\documentclass[../main.tex]{subfiles}
\begin{document}
\section*{第二部分提要}\label{section:part2-outline}

在\cref{chapter:ch06}中，我们研究如下初边值问题
\begin{align*}
\begin{cases}
\partial_t u = \Delta\phi(u) & \text{在 }\Omega\times(0,T),\\
u=0 & \text{在 }\partial\Omega\times(0,T),\\
u=u_0 & \text{在 }\Omega\times\{0\},
\end{cases}
\qquad(P)
\end{align*}
其中 $u_0$ 是 $\Omega\subseteq\mathbb{R}^N$（$N\geqslant3$）上的有限 Radon 测度，$\Omega$ 为开集且有界，函数 $\phi$ 非减，且在无穷远处至多为多项式增长。我们讨论该问题的适定性以及 Radon 测度值解的正则性结果，它们依赖于 $\phi$ 在无穷远处的性态，以及 $u_0$ 与合适的 Sobolev 容量相关的性质。

\cref{chapter:ch07}处理标量守恒律的 Cauchy 问题
\begin{align*}
\begin{cases}
\partial_t u + \partial_x[\phi(u)] = 0 & \text{在 }\mathbb{R}\times(0,T),\\
u=u_0 & \text{在 }\mathbb{R}\times\{0\},
\end{cases}
\qquad(CL1)
\end{align*}
其中 $u_0$ 是 $\mathbb{R}$ 上的正有限 Radon 测度，其奇异部分通常为 Dirac 质量的有限叠加，$\phi$ 在无穷远处为次线性。我们讨论该问题的适定性以及熵 Radon 测度值解的定性性质，它们依赖于 $u_0$ 与 $\phi$ 在无穷远处的性态。

最后，\cref{chapter:ch08}致力于研究如下不适定初边值问题的 Sobolev 正则化
\begin{align*}
\begin{cases}
\partial_t u = \nabla\cdot[\phi(\nabla u)] & \text{在 }\Omega\times(0,T),\\
u=0 & \text{在 }\partial\Omega\times(0,T),\\
u=u_0 & \text{在 }\Omega\times\{0\}.
\end{cases}
\qquad(FB)
\end{align*}
此处 $\Omega\subseteq\mathbb{R}^N$ 为开集且有界，$u_0$ 属于某个 Sobolev 空间，$\phi$ 在无穷远处满足合适的增长条件。通过研究正则化的粘性消失极限，证明 Young 测度值解的存在性。进而，依据 $\phi$ 的性质刻画极限 Young 测度，并研究解在大时间时的渐近性态。

\end{document}
